\scanpage{45}
\mspage{18}
On suppose, pour $i<j<k$
\[
(28)\qquad\Sigma_{ijk}=\Sigma_{ik}\cap\Sigma_{jk}\qquad
(\text{a priori }\Sigma_{ijk}\subset\Sigma_{ik}\cap\Sigma_{jk})
\]
\margin{Condition \textcircled{b}}
ce qui signifie aussi que $\forall z\in\Sigma_*$ au dessus d'un $k\in I$ (i.e. $z\in\Sigma_k$), l'hom. d'ens.
\[
\Sigma_*\longrightarrow I
\]
(\inserted{qui par \textcircled{a}} induit une application
\[
\Sigma_{*\leq z}\hookrightarrow I_{\leq k}
\]
\margin{\uline{injection}}
) induit
\[
(29)\qquad (\Sigma_*)_{\leq z}\hookrightarrow I_{\leq k}
\]
\margin{l'ordre de $(\Sigma_*)_z$ est \uline{induit} par celui de $I_{\leq k}$}
\margin{Ceci implique que $d,d'$ sont des drap. tels que $d''=d\cup d'$ soit un drapeau et que $d,d'$ y soient \unsure{cofinaux} i.e. $d,d'$ aient $=$ élément final $i$, alors $\Sigma_{d''}=\Sigma_d\cap\Sigma_{d'}$ dans $\Sigma_i$}

On suppose de plus que si $i,j\leq k$ n'ont pas de minorant \inserted{dans $I$}, alors
\[
\Sigma_{ik}\cap\Sigma_{jk}=\phi
\]
(si $i,j$ ont un minorant $l$, alors
\[
\Sigma_{ik}\cap\Sigma_{jk}\subset\ )
\]
\margin{Condition \textcircled{c} Gros}
\margin{Implique que si $d,d'$ sont deux drapeaux ayant $=$ élément final $k$, et si $d\cup d'$ n'est pas un drapeau, alors $\Sigma_d\cap\Sigma_{d'}=\phi$}
\margin{\textcircled{b} et \textcircled{c} impliquent que si $(d_\alpha)_\alpha$ est une famille \inserted{\unsure{finie}} de drapeaux de $I$, \deleted{\unsure{$\forall\alpha$}} \unsure{$\forall\alpha\in A$} \unsure{aient|sont} final $k$, donc les $\Sigma_{d_\alpha}\subset\Sigma_k$, alors $\bigcap\Sigma_{d_\alpha}$ est \uline{vide} si $\bigcup d_\alpha$ n'est pas un drapeau, et $=\Sigma_{\cup d_\alpha}$ sinon}

On suppose enfin que le système des $\Sigma_{ik}$ (pour $k$ fixé, \inserted{$i\in I_{<k}$}), qui sont des sous-espaces \deleted{bord} bord-ts de $\Sigma_k$, sont \uline{transversaux}, ce qui signifie (grâce : (30)\inserted{et 28,}) que si \unsure{on a} un drapeau $d=(i_{\unsure{0|1}},--i_n,k)$, alors les $\Sigma_{i_{\unsure{0|1}}k},\Sigma_{i_{\unsure{1|2}}k}\mathinner{\cdotp-}\Sigma_{i_n,k}$ dans $\Sigma_k$ sont des sous-ens \unsure{bordants|bord-ts} qui se coupent \uline{transversalement} le long de $\Sigma_d=\Sigma_{i_{\unsure{0|1}}-i_n,k}$.
\margin{Condition \textcircled{d}}
\margin{implique le $=$ dans les $\Sigma_d$\unsure{;|,} \unsure{on prend} les $\Sigma_{d'}\subset\Sigma_d$ (i.e. but pour $d'\supsetneqq d$, et $d$ cofinal dans $d'$) \unsure{ce qui} \unsure{les} bord-ts\inserted{$d'=d\cup\{i\}$} \unsure{qui sont} toutes pour dans $\Sigma_d$ \unsure{liss.|lisses}}
\margin{44}
